% !TeX root = ../../../thesis.tex

\subsubsection{\travisci}\label{sssec:travisci}
The final \CI{} platform which I will discuss is \travisci{}. This \CI{} system was launched in 2011 and can only be used in addition to an existing \github{} repository. \travisci{} build tasks can be configured in a similar fashion as GitLab CI, but the builds can exclusively be executed on their servers. Besides running builds after a commit has been pushed to the repository, it is also possible to schedule daily, weekly or monthly builds using cron jobs. Similar to \github{} Actions, open-source projects can be built at zero cost and a paid plan exists for private repositories \cite{travisanalysis}. It is not possible to create custom plugins, however \travisci{} already features built-in support for a variety of programming languages. In 2020, almost 1 million projects are being built using \travisci{} \cite{travis2020}.

\begin{figure}[htbp!]
	\centering
	\includegraphics[width=0.40\textwidth]{assets/images/travis-ci.pdf}
	\caption{Logo of \travisci{} (\url{https://travis-ci.com/})}
	\label{fig:travis-ci}
\end{figure}
